
\section{Entropy asymptotics}\label{app:entropy-asymptotics}
 This appendix proves the four large-\(k\) entropy estimates of the main
text: the single-output estimates of Corollary~\ref{cor:single-asymptotic}
and Section~\ref{sec:antisymmetric-postprocessing}, and the Bell-output
estimates of Sections~\ref{sec:bell-output}
and~\ref{sec:antisymmetric-postprocessing}.  All four proofs follow the
same scheme: we identify the relevant spectrum, write it as a
perturbation of a maximally mixed spectrum, and apply one exact entropy
identity together with one set of local bounds
(Propositions~\ref{prop:app-exact-expansion}
and~\ref{prop:app-local-bounds}).
 
Throughout, \(t\in(0,1)\), \(p\in(0,\infty)\), and an integer \(r\geq1\)
are fixed, \(\gamma:=(1-t)/t\), and \(k\to\infty\); in particular, we may
assume \(k\ge2r\) and \(k^2t>1\).  The constants \(C\), and the constants
implicit in \(O(\cdot)\), depend only on \(p\), \(r\), and \(t\); the value
of \(C\) may change from line to line.  We abbreviate \(D:=D_{k,r}\),
\(N:=N_{k,r}\), \(\Pi:=\Pi_{k,r}^{\asym}\), and \(\mc E:=\mc E_{k,r}\); for
\(r=1\), \(D=k\), \(N=1\), and \(\mc E\) is the identity channel.  A
probability vector \(\nu\) is \emph{written with multiplicities} if it
consists of values \(\nu_1,\ldots,\nu_n\) repeated
\(d_1,\ldots,d_n\geq1\) times, so that \(\sum_jd_j\nu_j=1\).
 
\subsection{R\'enyi entropy near a maximally mixed spectrum}
 
For \(x,y>-1\), define
\begin{equation}
    \psi_p(y):=
    \begin{cases}
        (1+y)^p-1-py, & p\neq1,\\
        (1+y)\log(1+y)-y, & p=1,
    \end{cases}
    \qquad
    F_p(x):=
    \begin{cases}
        \dfrac{\log(1+x)}{1-p}, & p\neq1,\\[6pt]
        -x, & p=1,
    \end{cases}
    \label{eq:app-psi-F}
\end{equation}
and put \(\kappa_p:=\frac12\psi_p''(0)\) and \(c_p:=F_p'(0)\).  Thus
\(\kappa_p=p(p-1)/2\) and \(c_p=1/(1-p)\) for \(p\neq1\), while
\(\kappa_1=1/2\) and \(c_1=-1\).  In all cases,
\begin{equation}
    c_p\kappa_p=-\frac p2.
    \label{eq:app-cp-kappa}
\end{equation}
 
\begin{proposition}[Exact expansion]
\label{prop:app-exact-expansion}
Let \(\nu\) be a probability vector written with multiplicities
\(d_1,\ldots,d_n\), put \(M:=\sum_jd_j\), and write
\(\nu_j=(1+y_j)/M\) with \(y_j>-1\).  Then \(\sum_jd_jy_j=0\) and
\begin{equation}
    S_p(\nu)
    =\log M+F_p\Bigl(\frac1M\sum_{j=1}^nd_j\,\psi_p(y_j)\Bigr).
    \label{eq:app-exact-expansion}
\end{equation}
\end{proposition}
 
\begin{proof}
Normalization gives \(\sum_jd_jy_j=0\). Put
\(X:=M^{-1}\sum_jd_j\psi_p(y_j)\). For \(p\neq1\),
\begin{equation}
    \sum_jd_j\nu_j^p
    =M^{-p}\sum_jd_j\bigl(1+py_j+\psi_p(y_j)\bigr)
    =M^{1-p}(1+X).
\end{equation}
In particular \(X>-1\), and taking logarithms proves the formula.
For \(p=1\), \(\psi_1(y)\geq0\), so \(X\geq0\), and
\begin{equation}
    S_1(\nu)
    =\log M-\frac1M\sum_jd_j(1+y_j)\log(1+y_j)
    =\log M-X.
\end{equation}
\end{proof}
 
\begin{proposition}[Local bounds]
\label{prop:app-local-bounds}
There are \(\eta_p\in(0,1/2]\) and \(C_p<\infty\) such that the following
hold.
\begin{enumerate}[label=\textup{(\roman*)}]
\item \(|\psi_p(y)-\kappa_py^2|\le C_p|y|^3\) for \(|y|\le\eta_p\).
\item \(|F_p(x)-c_px|\le C_px^2\) for \(|x|\le1/2\).
\item For every compact interval \(J\subset(-1,\infty)\), there is
\(L_J<\infty\) with \(|\psi_p(y)|\le L_J\) and
\(|\psi_p(y)-\psi_p(y')|\le L_J|y-y'|\) for all \(y,y'\in J\).
\end{enumerate}
\end{proposition}
 
\begin{proof}
Both functions are smooth on \((-1,\infty)\), with
\(\psi_p(0)=\psi_p'(0)=F_p(0)=0\),
\(\psi_p''(0)=2\kappa_p\), and \(F_p'(0)=c_p\).
Taylor's theorem gives (i)--(ii), and boundedness of \(\psi_p\) and
\(\psi_p'\) on \(J\) gives (iii).

For explicit constants, the elementary bounds
\begin{align}
    |\log(1+y)-y|&\leq2y^2,
    &|\log(1+y)-y+y^2/2|&\leq2|y|^3
       &&(|y|\leq1/2),\\
    |e^z-1-z-z^2/2|&\leq|z|^3
       &&&&(|z|\leq1)
\end{align}
give (i), for \(p\neq1\), with
\(\eta_p=\min\{1/2,1/(4p)\}\) and cubic constant
\(C_p^{\psi}=8p^3+3p^2+2p\). For \(p=1\), take
\(\eta_1=1/2\) and \(C_1^{\psi}=4\).
For (ii), take \(C_p^F=2/|1-p|\) when \(p\neq1\), and
\(C_1^F=0\). The common constant
\(C_p=\max\{C_p^{\psi},C_p^F\}\) proves the statement.
\end{proof}
 
\begin{proposition}[Uniform expansion]
\label{prop:app-uniform-expansion}
In the setting of Proposition~\ref{prop:app-exact-expansion}, put
\(\eta:=\max_j|y_j|\) and \(Q:=\frac1M\sum_jd_jy_j^2\).  After decreasing
\(\eta_p\) if necessary, \(\eta\le\eta_p\) implies
\begin{equation}
    \Bigl|S_p(\nu)-\log M+\frac p2\,Q\Bigr|\le C\,(\eta+Q)\,Q.
    \label{eq:app-uniform-expansion}
\end{equation}
\end{proposition}
 
\begin{proof}
Let \(X:=\frac1M\sum_jd_j\psi_p(y_j)\).  Since \(|y_j|^3\le\eta y_j^2\),
Proposition~\ref{prop:app-local-bounds}(i) gives
\(|X-\kappa_pQ|\le C_p\eta Q\), and hence \(|X|\le(|\kappa_p|+C_p)Q\).  Since \(Q\le\eta^2\), we may decrease
\(\eta_p\) so that \(|X|\le1/2\).  By \eqref{eq:app-cp-kappa},
\begin{equation}
    S_p(\nu)-\log M+\frac p2\,Q
    =\bigl(F_p(X)-c_pX\bigr)+c_p\bigl(X-\kappa_pQ\bigr),
\end{equation}
and Proposition~\ref{prop:app-local-bounds}(ii) bounds the absolute
value of the right-hand side by \(C_pX^2+|c_p|C_p\eta Q\le C(\eta+Q)Q\).
\end{proof}
 
\subsection{The single-output estimates}
 
\begin{proposition}[Single output]
\label{prop:app-single-entropy-asymptotics}
For every fixed \(t\in(0,1)\) and \(p\in(0,\infty)\),
\begin{equation}
    \min_{\sigma\in\mathscr K_{k,t}}S_p(\sigma)
    =\log k-\frac{2p(1-t)}{tk^2}+O_{p,t}(k^{-5/2}).
    \label{eq:app-single-entropy-asymptotics}
\end{equation}
\end{proposition}
 
\begin{proposition}[Single output after antisymmetric postprocessing]
\label{prop:app-antisymmetric-single-entropy-asymptotics}
For every fixed \(t\in(0,1)\), \(p\in(0,\infty)\), and integer \(r\geq1\),
\begin{equation}
    \min_{\rho\in\mathscr K_{k,t}}
        S_p\!\left(\mc E_{k,r}(\rho)\right)
    =
    \log\binom{k}{r}
    -\frac{2p(1-t)}{tr k^2}
    +O_{p,r,t}(k^{-5/2}).
    \label{eq:app-antisymmetric-single-entropy-asymptotics}
\end{equation}
\end{proposition}
 
Since \(\mc E_{k,1}\) is the identity channel,
Proposition~\ref{prop:app-single-entropy-asymptotics} is the case
\(r=1\) of
Proposition~\ref{prop:app-antisymmetric-single-entropy-asymptotics}.
We prove the latter for all \(r\geq1\) after two preparatory
propositions.  For a probability vector \(\lambda\in\mb R^k\), we write
\(\lambda=(\mathbf1+\varepsilon)/k\), where \(\mathbf1:=(1,\ldots,1)\), so
that \(\sum_i\varepsilon_i=0\).
 
\begin{proposition}[Shuffled spectrum]
\label{prop:app-shuffled-spectrum}
Let \(\rho\in\mc D_k\) have spectrum \(\lambda=(\mathbf1+\varepsilon)/k\).
Then the eigenvalues of \(\mc E(\rho)\), indexed by the \(r\)-subsets
\(I\subseteq[k]\), are \(\mu_I=(1+\delta_I)/D\) with
\(\delta_I:=\frac1r\sum_{i\in I}\varepsilon_i\).  Moreover,
\(\max_I|\delta_I|\le\|\varepsilon\|_\infty\) and
\begin{equation}
    \frac1D\sum_{|I|=r}\delta_I^2
    =\frac{k-r}{rk(k-1)}\,\|\varepsilon\|_2^2.
    \label{eq:app-quadratic-form}
\end{equation}
\end{proposition}
 
\begin{proof}
By Proposition~\ref{prop:eigenvalue-shuffling}, the eigenvalues of
\(\mc E(\rho)\) are
\(\mu_I=\frac1N\sum_{i\in I}\lambda_i=\frac1{kN}\bigl(r+\sum_{i\in I}\varepsilon_i\bigr)\),
and \(kN=rD\) by \eqref{eq:DN}.  This gives the formula for \(\mu_I\) and
the bound on \(\delta_I\).  For \eqref{eq:app-quadratic-form}, write
\(\varepsilon_I:=\sum_{i\in I}\varepsilon_i\).  Since
\(\sum_i\varepsilon_i=0\), we have
\(\varepsilon_I=-\sum_{j\notin I}\varepsilon_j\).  For \(i\neq j\),
exactly \(\binom{k-2}{r-1}\) of the \(r\)-subsets contain \(i\) but not
\(j\).  Hence
\begin{equation}
    \sum_{|I|=r}\varepsilon_I^2
    =-\sum_{|I|=r}\sum_{i\in I}\sum_{j\notin I}\varepsilon_i\varepsilon_j
    =-\binom{k-2}{r-1}\sum_{i\neq j}\varepsilon_i\varepsilon_j
    =\binom{k-2}{r-1}\|\varepsilon\|_2^2,
\end{equation}
and \(\binom{k-2}{r-1}/D=r(k-r)/(k(k-1))\).
\end{proof}
 
\begin{proposition}[Localization of \(\Lambda_{k,t}\)]
\label{prop:app-localization}
There are \(C\) and \(k_0\), depending only on \(t\), such that the
following hold for \(k\ge k_0\).
\begin{enumerate}[label=\textup{(\roman*)}]
\item Every \(\lambda\in\Lambda_{k,t}\) satisfies
\(\|\varepsilon\|_\infty\le Ck^{-1/2}\) and
\(\|\varepsilon\|_2^2\le\frac{4\gamma}k+Ck^{-3/2}\).
\item There is \(\lambda^*\in\Lambda_{k,t}\) with
\(\|\varepsilon^*\|_2^2\ge\frac{4\gamma}k-Ck^{-2}\).
\end{enumerate}
\end{proposition}
 
\begin{proof}
Put \(a=\sqrt t\), \(b=\sqrt{1-t}\), \(x=\sqrt u\), and
\(y=\sqrt{1-u}\). Since
\((ax+by)^2+(ay-bx)^2=1\) and \(ax+by\in[0,1]\),
\begin{equation}
    (x-a)^2+(y-b)^2=2(1-ax-by)\leq2c_t(u).
\end{equation}
Thus \(|bx+ay-2ab|\leq\sqrt{2c_t(u)}\). Using
\(u-t=(bx-ay)(bx+ay)\), we obtain
\begin{equation}
    |u-t|\leq\sqrt{c_t(u)}
       \left(2\sqrt{t(1-t)}+\sqrt{2c_t(u)}\right).
    \label{eq:app-ct-bound}
\end{equation}

For (i), let \(u\in\mathscr D_{k,t}\), and put
\(w_i=u_i-t\), \(c_i=c_t(u_i)\), and
\(\bar w=k^{-1}\sum_iw_i\). Since \(\sum_i c_i\leq1/k\),
\begin{equation}
    \max_i|w_i|\leq\frac2{\sqrt k},\qquad
    \sum_iw_i^2
       \leq\frac1k\left(2\sqrt{t(1-t)}+\sqrt{2/k}\right)^2
       \leq\frac{4t(1-t)}k+Ck^{-3/2}.
\end{equation}
The middle expression is at most \(4/k\) for \(k\geq2\), so
\(|\bar w|\leq2/k\). Now
\(\varepsilon_i=(w_i-\bar w)/(t+\bar w)\).
For \(k\geq4/t\), the denominator is at least \(t/2\); hence
\begin{equation}
    \|\varepsilon\|_\infty\leq Ck^{-1/2},\qquad
    \|\varepsilon\|_2^2
       \leq\frac{\sum_iw_i^2}{(t-2/k)^2}
       \leq\frac{4\gamma}k+Ck^{-3/2}.
\end{equation}

For (ii), put \(\sigma=(2k)^{-1/2}\), and assume
\(\sigma^2\leq\min\{t,1-t\}\). Define
\begin{equation}
    u_\pm=\left(a\sqrt{1-\sigma^2}\pm b\sigma\right)^2,
    \qquad u^*=(u_+,u_-,t,\ldots,t).
\end{equation}
Both square roots in
\(\sqrt{u_\pm}=a\sqrt{1-\sigma^2}\pm b\sigma\) and
\(\sqrt{1-u_\pm}=b\sqrt{1-\sigma^2}\mp a\sigma\)
are nonnegative. Therefore \(c_t(u_\pm)=\sigma^2\), and
\(u^*\in\mathscr D_{k,t}\). A direct calculation gives
\begin{equation}
    \bar w=\frac{1-2t}{k^2},\qquad
    \sum_i(w_i-\bar w)^2
    =\frac{4t(1-t)}k+
       \frac{(1-2t)^2/2-2t(1-t)}{k^2}
       -\frac{(1-2t)^2}{k^3}.
\end{equation}
Dividing by \((t+\bar w)^2\) yields
\(\|\varepsilon^*\|_2^2=4\gamma/k+O_t(k^{-2})\), proving (ii).
\end{proof}
 
\begin{proof}[Proof of Propositions~\ref{prop:app-single-entropy-asymptotics}
and~\ref{prop:app-antisymmetric-single-entropy-asymptotics}]
Write \(\lambda=(\mathbf1+\varepsilon)/k\in\Lambda_{k,t}\).
The shuffled-spectrum formula and localization give, uniformly in
\(\lambda\),
\begin{equation}
    \eta:=\max_I|\delta_I|=O_t(k^{-1/2}),\qquad
    Q:=\frac1D\sum_I\delta_I^2
       =\frac{k-r}{rk(k-1)}\|\varepsilon\|_2^2
       =O_{r,t}(k^{-2}).
\end{equation}
Proposition~\ref{prop:app-uniform-expansion} therefore gives
\begin{equation}
    S_p\!\left(\mc E_{k,r}(\rho)\right)
    =\log D-\frac{p(k-r)}{2rk(k-1)}\|\varepsilon\|_2^2
       +O_{p,r,t}(k^{-5/2}).
\end{equation}
The two bounds in Proposition~\ref{prop:app-localization} imply
\begin{equation}
    \sup_{\lambda\in\Lambda_{k,t}}\|\varepsilon\|_2^2
       =\frac{4\gamma}k+O_t(k^{-3/2}).
\end{equation}
Taking the minimum, using uniformity of the remainder and
\((k-r)/(k-1)=1+O_r(k^{-1})\), proves the postprocessing estimate.
Its case \(r=1\) is the single-output estimate.
\end{proof}
 
\subsection{The Bell-output estimates}
 
Recall that \(Z_{k,t}=r_{k,t}\psi_k^++(1-r_{k,t})I_{k^2}/k^2\) with
\begin{equation}
    r_{k,t}=\frac{k^2(1-t)}{k^4t-2k^2t+1}.
    \label{eq:app-rkt}
\end{equation}
Its spectrum \(\lambda^{\rm Bell}_{k,t}\) consists of
\(\alpha_{k,t}=r_{k,t}+(1-r_{k,t})/k^2\) with multiplicity \(1\) and
\(\beta_{k,t}=(1-r_{k,t})/k^2\) with multiplicity \(k^2-1\).  We put
\(\omega_{k,r,t}:=(\mc E_{k,r}\otimes\overline{\mc E}_{k,r})(Z_{k,t})\).
 
\begin{proposition}[Bell output]
\label{prop:app-bell-entropy-asymptotics}
For every fixed \(t\in(0,1)\) and \(p\in(0,\infty)\),
\begin{equation}
    S_p(\lambda^{\rm Bell}_{k,t})
    =2\log k-\frac{A_p(t)}{k^2}+O_{p,t}(k^{-4}),
    \label{eq:app-bell-entropy-asymptotics}
\end{equation}
where \(A_p(t):=\frac{t^{-p}-1-p(t^{-1}-1)}{p-1}\) for \(p\neq1\) and
\(A_1(t):=t^{-1}\log(t^{-1})-t^{-1}+1\).
\end{proposition}
 
\begin{proposition}[Bell output after antisymmetric postprocessing]
\label{prop:app-antisymmetric-bell-entropy-asymptotics}
For every fixed \(t\in(0,1)\), \(p\in(0,\infty)\), and integer \(r\geq1\),
\begin{equation}
    S_p(\omega_{k,r,t})
       =
       2\log\binom{k}{r}
       -\frac{B_{p,r}(\gamma)}{k^2}
       +O_{p,r,t}(k^{-3}),
    \label{eq:app-antisymmetric-bell-entropy-asymptotics}
\end{equation}
where
\begin{equation}
    B_{p,r}(\gamma)
       :=
       \begin{cases}
       \displaystyle
       \frac{p\gamma-r^2\bigl[(1+\gamma/r^2)^p-1\bigr]}{1-p},
           & p\neq1,\\[8pt]
       \displaystyle
       (r^2+\gamma)\log(1+\gamma/r^2)-\gamma,
           & p=1.
       \end{cases}
\end{equation}
\end{proposition}
 
Since \(\mc E_{k,1}\) is the identity channel,
\(\omega_{k,1,t}=Z_{k,t}\).  Since moreover \(1+\gamma=t^{-1}\), we have
\(B_{p,1}(\gamma)=A_p(t)\).  Thus
Proposition~\ref{prop:app-bell-entropy-asymptotics} is the case \(r=1\)
of Proposition~\ref{prop:app-antisymmetric-bell-entropy-asymptotics},
with a sharper error term.  In both cases, \eqref{eq:app-psi-F} gives
\begin{equation}
    c_p\,r^2\,\psi_p(\gamma/r^2)=-B_{p,r}(\gamma).
    \label{eq:app-B-identity}
\end{equation}
For \(r\geq2\), the spectrum of \(\omega_{k,r,t}\) requires the
following four propositions; for \(r=1\), it is the explicit spectrum
\(\lambda^{\rm Bell}_{k,t}\).
 
For \(0\le m\le r\), write \(\Pi_m:=\Pi_{k,m}^{\asym}\) and
\(\mc A_m:=\mc A_{k,m}^{\asym}\), with \(\mc A_0:=\mb C\) and
\(\Pi_0:=1\); thus \(\Pi=\Pi_r\).  We identify \(\mc B(\mc A_m)\) with
the operators \(Y\) on \((\mb C^k)^{\otimes m}\) satisfying
\(Y=\Pi_mY\Pi_m\), equipped with the Hilbert--Schmidt inner product
\(\langle Y,Y'\rangle=\Tr(Y^*Y')\).  Let \(F_{jl}\) be the swap of the
\(j\)th and \(l\)th tensor factors.  For an operator \(Y\) on
\((\mb C^k)^{\otimes m}\), \((\Tr_jY)\otimes I_{(j)}\) denotes the
operator obtained by tracing out the \(j\)th factor and then placing
\(I_k\) on that factor, keeping all other factors in place.
 
Each \(\pi\in\mathfrak S_m\) can be written uniquely as \(\pi=(l\ m)\sigma\)
with \(1\le l\le m\) and \(\sigma\) fixing \(m\), where \((m\ m)\) denotes
the identity.  Taking the signed average over \(\pi\) gives
\begin{equation}
    \Pi_m
    =\frac1m\Bigl(I-\sum_{l=1}^{m-1}F_{lm}\Bigr)(\Pi_{m-1}\otimes I_k)
    =\frac1m(\Pi_{m-1}\otimes I_k)\Bigl(I-\sum_{l=1}^{m-1}F_{lm}\Bigr).
    \label{eq:app-antisym-recursion}
\end{equation}
In particular,
\(\Pi_m=(\Pi_{m-1}\otimes I_k)\Pi_m=\Pi_m(\Pi_{m-1}\otimes I_k)\).
Moreover, \(\Pi_mF_{jl}=F_{jl}\Pi_m=-\Pi_m\) for \(j\neq l\).
 
For \(1\le m\le r\), the one-particle reduction
\begin{equation}
    R_m:\mc B(\mc A_m)\to\mc B(\mc A_{m-1}),
    \qquad
    R_m(Y):=\Tr_mY,
\end{equation}
is well defined: if \(Y=\Pi_mY\Pi_m\), then
\(Y=(\Pi_{m-1}\otimes I_k)Y(\Pi_{m-1}\otimes I_k)\), and hence
\(\Tr_mY=\Pi_{m-1}(\Tr_mY)\Pi_{m-1}\).  Its Hilbert--Schmidt adjoint is
the one-particle lift \(R_m^*(Z)=\Pi_m(Z\otimes I_k)\Pi_m\), because
\(\Tr\bigl(Z^*\Tr_mY\bigr)=\Tr\bigl((Z^*\otimes I_k)Y\bigr)
=\Tr\bigl((\Pi_m(Z\otimes I_k)\Pi_m)^*Y\bigr)\).  We write
\begin{equation}
    G_m:=R_m^*R_m\ \text{on }\mc B(\mc A_m),
    \qquad
    H_m:=R_mR_m^*\ \text{on }\mc B(\mc A_{m-1}).
\end{equation}
Both are positive semidefinite, and they have the same nonzero
eigenvalues, with multiplicities.  For \(a,b\in[k]\), the collective
one-body operators on \(\mc A_m\) are
\begin{equation}
 L_{ab}^{(m)}
 :=\Pi_m\Bigl(\sum_{j=1}^mE_{ab}^{(j)}\Bigr)\Pi_m,
 \qquad
 E_{ab}:=|a\rangle\langle b|.
 \label{eq:app-Lab-m}
\end{equation}
By \eqref{eq:E-definition}, \(\mc E(E_{ab})=N^{-1}L_{ab}^{(r)}\).
 
\begin{proposition}[One-body twirl]\label{prop:app-one-body-twirl}
For \(1\le m\le r\) and \(Y\in\mc B(\mc A_m)\),
\begin{equation}
    \sum_{a,b=1}^kL_{ab}^{(m)}\,Y\,L_{ba}^{(m)}
    =m^2\,G_m(Y).
\end{equation}
\end{proposition}
 
\begin{proof}
For a single factor, \(\sum_{a,b}E_{ab}XE_{ba}=\Tr(X)I_k\).  Applied to
the \(j\)th factor, this gives
\(\sum_{a,b}E_{ab}^{(j)}XE_{ba}^{(j)}=(\Tr_jX)\otimes I_{(j)}\).  Now let
\(j\neq l\).  Then \(E_{ba}^{(l)}=F_{jl}E_{ba}^{(j)}F_{jl}\), and
\(YF_{jl}=Y\Pi_mF_{jl}=-Y\).  Hence
\begin{equation}
    \sum_{a,b}E_{ab}^{(j)}\,Y\,E_{ba}^{(l)}
    =\bigl((\Tr_j(YF_{jl}))\otimes I_{(j)}\bigr)F_{jl}
    =-\bigl((\Tr_jY)\otimes I_{(j)}\bigr)F_{jl}.
\end{equation}
Since \(Y=\Pi_mY\Pi_m\), we have
\(\sum_{a,b}L_{ab}^{(m)}YL_{ba}^{(m)}
=\Pi_m\bigl(\sum_{j,l}\sum_{a,b}E_{ab}^{(j)}YE_{ba}^{(l)}\bigr)\Pi_m\).
Inserting the two identities above and using \(F_{jl}\Pi_m=-\Pi_m\), we
obtain
\begin{equation}
    \sum_{a,b}L_{ab}^{(m)}YL_{ba}^{(m)}
    =m\sum_{j=1}^m\Pi_m\bigl((\Tr_jY)\otimes I_{(j)}\bigr)\Pi_m.
\end{equation}
Finally, conjugation by \(F_{jm}\) leaves \(Y\) invariant and maps
\((\Tr_jY)\otimes I_{(j)}\) to \((\Tr_mY)\otimes I_{(m)}\).  Since
\(\Pi_mF_{jm}=-\Pi_m\), each summand equals
\(\Pi_m\bigl((\Tr_mY)\otimes I_k\bigr)\Pi_m=G_m(Y)\).
\end{proof}
 
\begin{proposition}[Reduction after lifting]\label{prop:app-lift-reduce}
We have \(H_1=k\) on \(\mc B(\mc A_0)=\mb C\), and for \(2\le m\le r\),
\begin{equation}
    H_m=\frac1{m^2}\Bigl((k-2m+2)\,\mathrm{id}+(m-1)^2\,G_{m-1}\Bigr)
    \qquad\text{on }\mc B(\mc A_{m-1}).
    \label{eq:app-lift-reduce}
\end{equation}
\end{proposition}
 
\begin{proof}
For \(m=1\), \(H_1(z)=\Tr(zI_k)=kz\).  Let \(m\geq2\) and
\(Z\in\mc B(\mc A_{m-1})\), and put \(\widehat Z:=Z\otimes I_k\).  By
\eqref{eq:app-antisym-recursion},
\begin{equation}
\begin{aligned}
    H_m(Z)&=\Tr_m\bigl(\Pi_m\widehat Z\Pi_m\bigr)\\
    &=\frac1{m^2}\,\Pi_{m-1}\Tr_m\Bigl[
       \Bigl(I-\sum_jF_{jm}\Bigr)\widehat Z\Bigl(I-\sum_lF_{lm}\Bigr)
     \Bigr]\Pi_{m-1},
\end{aligned}
\end{equation}
where \(j,l\) range over \([m-1]\).  We evaluate the terms.  First,
\(\Tr_m\widehat Z=kZ\).  Second,
\(\Tr_m(\widehat ZF_{lm})=Z\Tr_m(F_{lm})=Z\), and likewise
\(\Tr_m(F_{jm}\widehat Z)=Z\); these terms contribute \(-2(m-1)Z\).
Third, \(\Tr_m(F_{jm}\widehat ZF_{jm})=(\Tr_jZ)\otimes I_{(j)}\).
Fourth, let \(j\neq l\).  Then \(F_{lm}=F_{jl}F_{jm}F_{jl}\) and
\(\widehat ZF_{jl}=-\widehat Z\), so
\(F_{jm}\widehat ZF_{lm}=-F_{jm}\widehat ZF_{jm}F_{jl}\), and hence
\(\Tr_m(F_{jm}\widehat ZF_{lm})=-\bigl((\Tr_jZ)\otimes I_{(j)}\bigr)F_{jl}\).
After compression by \(\Pi_{m-1}\), using \(F_{jl}\Pi_{m-1}=-\Pi_{m-1}\)
and the last step of the proof of Proposition~\ref{prop:app-one-body-twirl},
each of the \((m-1)^2\) terms of the third and fourth types equals
\(G_{m-1}(Z)\).  Since \(\Pi_{m-1}Z\Pi_{m-1}=Z\), we obtain
\(H_m(Z)=\frac1{m^2}\bigl((k-2m+2)Z+(m-1)^2G_{m-1}(Z)\bigr)\).
\end{proof}
 
\begin{proposition}\label{prop:app-spectrum-G}
Let \(k\ge2r\), and for \(0\le m\le r\) put
\begin{equation}
    \lambda_m^{(r)}:=\frac{(r-m)(k-r-m+1)}{r^2},
    \qquad
    d_m:=\binom km^2-\binom k{m-1}^2,
\end{equation}
with \(\binom k{-1}:=0\).  Then \(G_r\) has eigenvalues
\(\lambda_m^{(r)}\) with multiplicities \(d_m\), \(0\le m\le r\).
\end{proposition}
 
\begin{proof}
We argue by induction on \(r\).  For \(r=1\), \(H_1=k\).  Hence \(G_1\)
has the eigenvalue \(k=\lambda_0^{(1)}\) with multiplicity \(1=d_0\), and
the eigenvalue \(0=\lambda_1^{(1)}\) with multiplicity \(k^2-1=d_1\).
Let \(r\geq2\), and assume the claim for \(r-1\).  By
Proposition~\ref{prop:app-lift-reduce}, \(H_r\) has the eigenvalues
\(\frac1{r^2}\bigl(k-2r+2+(r-1)^2\lambda_m^{(r-1)}\bigr)\) with
multiplicities \(d_m\), \(0\le m\le r-1\).  The identity
\((r-1-m)(k-r-m+2)+k-2r+2=(r-m)(k-r-m+1)\) shows that these eigenvalues
are \(\lambda_m^{(r)}\).  They are at least
\(\lambda_{r-1}^{(r)}=(k-2r+2)/r^2>0\).  Hence the nonzero eigenvalues
of \(G_r\) are \(\lambda_m^{(r)}\), \(0\le m\le r-1\), with multiplicities
\(d_m\).  The eigenvalue \(0=\lambda_r^{(r)}\) has multiplicity
\(\dim\mc B(\mc A_r)-\sum_{m=0}^{r-1}d_m=\binom kr^2-\binom k{r-1}^2=d_r\).
\end{proof}
 
\begin{proposition}[Spectrum of the post-processed Bell state]
\label{prop:app-spectrum-W}
Let \(k\ge2r\).  The state \(W_{k,r}\) of \eqref{eq:W-direct} has eigenvalues
\begin{equation}
    w_m:=\frac{(r-m)(k-r-m+1)}{kN^2},
    \qquad 0\le m\le r,
\end{equation}
with multiplicities \(d_m\).
\end{proposition}
 
\begin{proof}
The matrix entries of \(L_{ab}^{(r)}\) are real, so by
\eqref{eq:W-direct},
\begin{equation}
    W_{k,r}=\frac1{kN^2}\sum_{a,b=1}^kL_{ab}^{(r)}\otimes L_{ab}^{(r)}.
\end{equation}
Let \(\mathrm{vec}:\mc B(\mc A_r)\to\mc A_r\otimes\mc A_r\) be the unitary
with \(\mathrm{vec}\bigl(|e_I^{\asym}\rangle\langle e_J^{\asym}|\bigr)
=|e_I^{\asym}\rangle\otimes|e_J^{\asym}\rangle\).  Then
\((X_1\otimes X_2)\mathrm{vec}(Y)=\mathrm{vec}(X_1YX_2^T)\), with the
transpose taken in the real Slater basis.  The matrix entries of
\(L_{ab}^{(r)}\) in this basis are real, and
\((L_{ab}^{(r)})^*=L_{ba}^{(r)}\), so \((L_{ab}^{(r)})^T=L_{ba}^{(r)}\).
Therefore, by Proposition~\ref{prop:app-one-body-twirl},
\begin{equation}
    W_{k,r}\,\mathrm{vec}(Y)
    =\frac1{kN^2}\,\mathrm{vec}\Bigl(\sum_{a,b}L_{ab}^{(r)}YL_{ba}^{(r)}\Bigr)
    =\frac{r^2}{kN^2}\,\mathrm{vec}\bigl(G_r(Y)\bigr).
\end{equation}
Thus \(W_{k,r}\) is unitarily equivalent to \(\frac{r^2}{kN^2}G_r\), and
Proposition~\ref{prop:app-spectrum-G} gives the claim.
\end{proof}
 
\begin{proof}[Proof of Propositions~\ref{prop:app-bell-entropy-asymptotics}
and~\ref{prop:app-antisymmetric-bell-entropy-asymptotics}]
Put \(\varrho=r_{k,t}\). Since \(\mc E(I_k/k)=\Pi/D\), the
spectrum of \(\omega_{k,r,t}\), with multiplicities \(d_m\), is
\(\nu_m=(1+y_m)/D^2\), where
\begin{equation}
    y_m=\varrho(D^2w_m-1)
       =\varrho\left(\frac{k(r-m)(k-r-m+1)}{r^2}-1\right).
\end{equation}
Here \(0<\varrho<1\) for all sufficiently large \(k\), so \(y_m>-1\).
The exact expansion gives
\begin{equation}
    S_p(\omega_{k,r,t})=2\log D+F_p(X),\qquad
    X:=\frac1{D^2}\sum_{m=0}^rd_m\psi_p(y_m).
    \label{eq:app-bell-X}
\end{equation}

We have \(\varrho=\gamma/k^2+O_t(k^{-4})\), and all \(y_m\) lie in
\([-1/2,2\gamma]\) for large \(k\). The term \(m=r\) has
\(y_r=-\varrho\), so its contribution to \(X\) is \(O(k^{-4})\).
For \(m=r-1\),
\begin{align}
    y_{r-1}&=\frac\gamma{r^2}+O_{r,t}(k^{-1}),\\
    \frac{d_{r-1}}{D^2}
    &=\left(\frac r{k-r+1}\right)^2
      \left(1-\left(\frac{r-1}{k-r+2}\right)^2\right)
      =\frac{r^2}{k^2}+O_r(k^{-3}).
\end{align}
The local Lipschitz bound on \(\psi_p\) therefore makes this
contribution \(r^2\psi_p(\gamma/r^2)/k^2+O(k^{-3})\).
For \(m\leq r-2\), the successive binomial ratios are at most
\(2r/k\), and hence
\(d_m/D^2\leq\binom km^2/\binom kr^2\leq(2r/k)^4\).
These finitely many terms contribute \(O(k^{-4})\). Thus
\begin{equation}
    X=\frac{r^2}{k^2}\psi_p(\gamma/r^2)+O_{p,r,t}(k^{-3}).
    \label{eq:app-bell-X-estimate}
\end{equation}
Since \(F_p(X)=c_pX+O_p(X^2)\), identity
\eqref{eq:app-B-identity} proves the postprocessing estimate.

For \(r=1\), the exceptional multiplicity is exactly
\(d_0/D^2=1/k^2\), and \(y_0=\gamma+O_t(k^{-2})\).
There are no terms with \(m\leq r-2\), so the remainder in
\eqref{eq:app-bell-X-estimate} improves to \(O_{p,t}(k^{-4})\).
Using \(D=k\) and \(B_{p,1}(\gamma)=A_p(t)\) proves the Bell estimate.
\end{proof}

